\chapter{Données alternatives et données textuelles}\label{ch:rs:alternative-and-text-data}

Un fournisseur propose cinq ans de données de dépenses par carte portant sur quatre cents détaillants, pour des frais de plusieurs centaines de milliers de dollars par
an. Le fichier d'essai arrive accompagné d'une note~: le panel a changé de fournisseur de cartes la troisième année. L'historique paraît excellent. Les
questions sont celles du chapitre~: qu'est-ce exactement que ce jeu de données, quelle partie de son historique était connue à l'époque, ce qu'il ajoute à
ce que la société possède déjà, s'il est légal de l'utiliser, et ce qu'il vaut. Froot, Kang, Ozik et Sadka ont construit des indicateurs approchés en temps réel des
ventes de détail à partir de sources comprenant environ 50 millions d'appareils mobiles, et ont constaté que leur mesure intra-trimestrielle expliquait
la croissance des ventes et les \omterm{def:rs:fundamental-analyst-and-event-features:sue}{surprises de résultats}, avec des rendements anormaux moyens de 3.4~\% à l'annonce~; Zhu (2019) a montré que
l'introduction de données de transactions de consommation et de données satellitaires accroissait le contenu informatif des cours. De telles données sont exploitables. L'achat est
un projet de recherche à part entière, avec un cadre d'évaluation (\texttt{firm.vendoreval}), un essai synthétique dont la vérité est introduite à dessein, et une décision
à la clé.

\section{Ce que sont les jeux de données}

\begin{definition}[Données alternatives, dérive du panel]\label{def:rs:alternative-and-text-data:altdata}
Les \emph{données alternatives}\index{données alternatives} sont des données portant sur des entreprises, des économies ou des marchés, provenant de sources autres que les cours, les dépôts réglementaires et les statistiques officielles~:
transactions par carte et par banque, tickets de caisse, trafic et prix sur le web, usage d'applications, géolocalisation, images satellitaires
et aériennes, registres de fret et de vols, offres d'emploi, textes. Un jeu de données construit à partir d'un panel (de porteurs de cartes, d'appareils, de magasins) subit une
\emph{dérive du panel}\index{dérive du panel} lorsque la composition du panel, ou son lien avec la grandeur mesurée, évolue dans le temps~:
changement de fournisseur, attrition, nouvelle source, repondération.
\end{definition}

Tout jeu de données alternatif est un processus de mesure doté d'une histoire propre, et trois questions précèdent tout backtest. Que
mesure-t-on, par qui, et avec quel consentement (la provenance)~? Quand chaque valeur historique a-t-elle été délivrée pour la première fois (le
\omterm{def:rs:point-in-time-data-and-the-biases:bitemporal}{temps de connaissance} du chapitre~3)~? Comment le panel a-t-il évolué (sa dérive)~? Les deux dernières questions sont généralement liées~: un fournisseur lance son produit, puis reconstitue plusieurs années
d'historique à partir du panel actuel afin de vendre un backtest plus long, et cette reconstitution est ajustée, sciemment ou non, aux résultats déjà
connus au moment de sa construction. C'est le \omterm{def:rs:point-in-time-data-and-the-biases:vintage}{biais de remplissage rétroactif} (chapitre~3) sous sa forme commerciale.

\section{Évaluer un fournisseur}

\begin{definition}[Essai de données, coefficient d'information incrémental]\label{def:rs:alternative-and-text-data:trial}
Un \emph{essai de données}\index{essai de données} est une évaluation limitée dans le temps d'un jeu de données, menée sous licence d'essai, avant l'achat. Le
\emph{coefficient d'information incrémental}\index{coefficient d'information incrémental} d'une nouvelle variable explicative est le \omterm{def:rs:anatomy-of-a-predictor:ic}{coefficient d'information}
de son résidu, obtenu après régression de la variable, coupe transversale par coupe transversale, sur les variables explicatives que la société possède déjà.
\end{definition}

L'essai du chapitre est synthétique et son récit est introduit à dessein, si bien que chaque conclusion peut être vérifiée. Dans \texttt{firm.synthmkt}, les 400 valeurs
cotées pendant les dix années sont les détaillants~; la vraie \omterm{def:rs:fundamental-analyst-and-event-features:sue}{surprise de résultats} de chaque trimestre est annoncée un nombre aléatoire de jours après la
fin du trimestre, avec un saut de cours dans son sens. La valeur du panel fournie par le fournisseur pour chaque détaillant et chaque trimestre est livrée dix jours de bourse
après la fin du trimestre. Son historique de cinq ans (années 6 à 10 du marché) cache trois éléments~: le fournisseur est entré en service lorsqu'il a changé
de fournisseur de panel, au début de la troisième année de l'historique~; les deux premières années ont été reconstituées au lancement à partir du panel du nouveau fournisseur,
repondéré pour correspondre aux résultats publiés par les détaillants~; et le panel en service couvre moins de valeurs, avec un décalage de niveau et davantage de bruit. La société
possède déjà un signal qui porte une partie de la même information. L'essai ne retient que les trimestres dont l'annonce tombe après la
livraison, soit 5\,980 lignes, et \texttt{firm.vendoreval} mesure ce qu'il peut observer.

\textbf{Couverture et historique.} Les années reconstituées comptent 338 lignes par trimestre, les années en service 273~: l'échantillon s'est réduit au lancement du produit,
à l'inverse d'un panel en croissance, et c'est un indice. Sur l'ensemble des lignes, 39.6~\% ont été délivrées pour la première fois plus d'un trimestre après leur période (la
date de lancement du fournisseur, que l'essai doit demander, permet de les dater).

\textbf{Le lien avec la vérité.} Régression de la valeur du panel sur la surprise déclarée, année par année
(\cref{fig:rs:alternative-and-text-data:drift})~: pente de 0.92 et 0.87 les deux premières années, puis 0.37, 0.39 et 0.35~; corrélation
de 0.85 et 0.84, puis 0.35, 0.35 et 0.32~; ordonnée à l'origine de 0.02 avant, puis de 0.12 à 0.17 après. Une CUSUM des résidus, ordonnés dans le temps (chapitre~13),
place la rupture la plus vraisemblable au début de la troisième année, avec une statistique de 3.05 (la valeur critique à 5~\% est d'environ 1.36). Le
changement de fournisseur que le vendeur a mentionné est aussi le point à partir duquel l'historique cesse d'être une reconstitution.

\begin{omfigure}
\begin{tikzpicture}
\begin{axis}[width=10cm,height=6cm,xlabel={surprise de résultats déclarée (standardisée)},ylabel={valeur du panel du fournisseur},
  tick label style={font=\small},label style={font=\small},xmin=-3.5,xmax=3.5,ymin=-4,ymax=4.5,grid=major,grid style={gray!25},
  legend columns=3,
  legend style={font=\footnotesize,at={(0.5,-0.28)},anchor=north,draw=none,fill=none,
  /tikz/every even column/.append style={column sep=8pt}},table/col sep=comma]
\addplot[only marks,mark=*,mark size=0.7pt,gray!70] table[x=s,y=m]{figdata/research/12-alternative-and-text-data/live.csv};
\addplot[only marks,mark=*,mark size=0.7pt,omDef] table[x=s,y=m]{figdata/research/12-alternative-and-text-data/backfill.csv};
\addplot[omThm,very thick] table[x=s,y=m]{figdata/research/12-alternative-and-text-data/fitline.csv};
\legend{années en service 3--5,années reconstituées 1--2,ajustement des années en service}
\end{axis}
\end{tikzpicture}

\omcaption{La valeur du panel du fournisseur en fonction de la \omterm{def:rs:fundamental-analyst-and-event-features:sue}{surprise de résultats} déclarée par les détaillants~: un échantillon de 600 lignes pour chaque régime. L'historique reconstitué
épouse étroitement le résultat auquel il a été ajusté~; le panel en service est décalé et plus bruité. Données~: l'essai synthétique du chapitre sur
\texttt{firm.synthmkt}.}\label{fig:rs:alternative-and-text-data:drift}
\end{omfigure}

\textbf{Le \omterm{def:rs:anatomy-of-a-predictor:ic}{coefficient d'information} par année} avec le rendement entre la livraison et l'annonce~: 0.38 et 0.37 les années reconstituées, 0.12,
0.14 et 0.15 les années en service (\cref{fig:rs:alternative-and-text-data:ic}). Un backtest sur l'ensemble de l'historique rapporterait la moyenne des
deux régimes et ne décrirait ni l'un ni l'autre.

\textbf{Ce qui est nouveau.} Pendant les années en service, le signal propre à la société a un IC de 0.20, supérieur à celui du fournisseur, qui est de 0.13. Résiduée sur celui-ci trimestre par
trimestre, la variable du fournisseur conserve un \omterm{def:rs:alternative-and-text-data:trial}{coefficient d'information incrémental} de 0.10, avec une statistique $t$ de 4.0 sur 12 trimestres~: le panel apporte une information
qui manque à la société, moindre que ce que suggère son IC brut.

\begin{omfigure}
\begin{tikzpicture}
\begin{axis}[width=10cm,height=5.4cm,bar width=14pt,xlabel={année de l'historique du fournisseur},ylabel={IC de rang},tick label style={font=\small},
  label style={font=\small},xtick={1,2,3,4,5},xmin=0.4,xmax=5.6,ymin=0,ymax=0.45,grid=major,grid style={gray!25},
  legend columns=1,
  legend style={font=\footnotesize,at={(0.5,-0.28)},anchor=north,draw=none,fill=none,
  /tikz/every even column/.append style={column sep=8pt}},table/col sep=comma]
\addplot[ybar,area legend,fill=omDef!55,draw=omDef] table[x=year,y=ic]{figdata/research/12-alternative-and-text-data/ic.csv};
\addplot[omThm,thick,dashed,sharp plot,mark=none] table[x=year,y=inc]{figdata/research/12-alternative-and-text-data/icref.csv};
\addplot[black!60,thick,dotted,sharp plot,mark=none] table[x=year,y=own]{figdata/research/12-alternative-and-text-data/icref.csv};
\legend{panel du fournisseur,coefficient d'information incrémental (en service),signal propre de la société (en service)}
\end{axis}
\end{tikzpicture}

\omcaption{Rang IC de la valeur du panel du fournisseur avec le rendement entre la livraison et l'annonce, par année de son historique~; pour les années en service,
l'IC du signal que possède déjà la société et l'IC incrémental du fournisseur par rapport à ce signal. Données~: l'essai synthétique du
chapitre.}\label{fig:rs:alternative-and-text-data:ic}
\end{omfigure}

Les équivalents sur données réelles de chacune de ces étapes sont bien connus de quiconque a mené des essais~: tableaux de couverture avec des trous lorsqu'un marchand quitte le
panel, historiques dont le fournisseur ne peut pas fournir la date de première livraison, tickers rapprochés du \omterm{def:rs:universe-symbology-and-corporate-actions:master}{référentiel des titres} actuel (chapitre~4), données révisées
sans préavis. Le cadre d'évaluation pose les mêmes questions à tout fichier.

\section{Le texte comme donnée}

\begin{definition}[Méthode par dictionnaire, score de sentiment, plongement de document]\label{def:rs:alternative-and-text-data:text}
Une \emph{méthode par dictionnaire}\index{méthode par dictionnaire} attribue un score à un document en comptant les mots qu'il contient à partir de listes de mots fixées à l'avance. Un
\emph{score de sentiment}\index{score de sentiment} est une mesure de ce type, par exemple la part de mots négatifs, ou la différence entre les mots négatifs
et les mots positifs, le tout rapporté à l'ensemble des mots. Un \emph{plongement de document}\index{plongement de document} (embedding) associe à un document un vecteur de nombres réels (comptages
de mots pondérés par leur rareté, ou sortie d'un modèle de langage entraîné) de sorte que des documents au contenu voisin soient représentés par des vecteurs proches.
\end{definition}

Tetlock (2007) a mesuré le pessimisme d'une chronique quotidienne du \emph{Wall Street Journal} et constaté qu'un pessimisme médiatique élevé prédit une pression à la baisse sur les cours, suivie d'un retour aux fondamentaux, et
qu'un pessimisme anormalement élevé ou faible prédit un volume de transactions important.
Loughran et McDonald (2011) ont montré que des listes de mots construites pour d'autres disciplines classent à tort des mots courants dans un texte financier~: dans les 10-K
de 1994 à 2008, près des trois quarts des mots que le dictionnaire de Harvard, largement utilisé, classe comme négatifs ne le sont pas dans un
contexte financier (liability, tax, cost, capital), et ils ont donc construit des listes propres à la finance.

Le corpus du chapitre met en évidence ce problème. Deux mille documents synthétiques de 400 mots ont un ton $\tau$ qui augmente le taux de huit
mots positifs et diminue celui de huit mots négatifs~; huit mots financiers qu'une liste générale classe comme négatifs sont fréquents et sans lien avec
le ton, et représentent 77~\% des occurrences négatives de la liste générale. La part de mots négatifs est corrélée à $-\tau$ à 0.74 avec la liste générale, à
0.87 avec la liste financière, et à 0.95 pour la différence entre mots négatifs et mots positifs avec les listes financières. Les nombres sont ceux du corpus~; la leçon est celle de Loughran
et McDonald~: un dictionnaire est un instrument de mesure, calibré pour un domaine.

Les plongements remplacent les listes construites à la main par des représentations apprises, et déplacent les problèmes au lieu de les résoudre. Leur corpus d'entraînement a
une date, et un modèle entraîné sur des textes rédigés après la date de décision du backtest connaît la fin de l'histoire~; son vocabulaire et sa notion de ce qui est semblable
dérive d'une version à l'autre~; et l'application qui relie le plongement au rendement est un modèle ajusté de plus, avec les risques de surajustement du
chapitre~20. La règle du point-in-time s'applique au modèle comme aux données~: le modèle utilisé à une date de décision est celui qui existait alors.

\section{Droit et éthique}

\begin{definition}[Information privilégiée]\label{def:rs:alternative-and-text-data:mnpi}
L'\emph{information privilégiée}\index{information privilégiée} (material non-public information, MNPI) est une information sur un titre ou sur son émetteur qui n'a pas
été rendue publique et qui est significative~: selon les termes de la SEC, un élément est significatif s'il existe une probabilité substantielle qu'une personne raisonnable
le juge important. Opérer sur une information privilégiée obtenue en violation d'un devoir constitue un délit d'initié.
\end{definition}

Les \omterm{def:rs:alternative-and-text-data:altdata}{données alternatives} se situent près de cette limite par construction~: elles ont de la valeur parce qu'elles ne figurent pas dans les dépôts réglementaires. Les questions juridiques portent sur la manière dont les
données ont été obtenues et sur le droit du vendeur de les céder, et la diligence de l'acheteur fait partie de sa propre défense.

\begin{dated}{2026-09}{Données alternatives devant les autorités de contrôle et les tribunaux}\label{dat:rs:alternative-and-text-data:law}
En septembre 2021, la SEC a réglé des poursuites pour fraude sur titres contre App Annie, un fournisseur de \omterm{def:rs:alternative-and-text-data:altdata}{données alternatives} pour l'industrie des applications mobiles, ainsi que
son cofondateur, qui a accepté de payer plus de 10~millions de dollars~: c'est la première action d'application de la SEC visant un fournisseur de \omterm{def:rs:alternative-and-text-data:altdata}{données alternatives}
pour fraude sur titres. L'ordonnance a établi que, tout en assurant aux entreprises que leurs données confidentielles sur les applications seraient agrégées et anonymisées, App Annie a utilisé
des données non agrégées et non anonymisées de la fin 2014 au milieu de 2018 pour modifier ses estimations et les rendre plus précieuses pour les sociétés de trading, et
qu'elle a présenté cela de façon trompeuse à ces sociétés. En avril 2022, la cour d'appel des États-Unis pour le neuvième circuit, dans \emph{hiQ Labs v.\ LinkedIn}, a de nouveau
confirmé une injonction préliminaire permettant à une entreprise d'analyse de données d'accéder aux profils LinkedIn accessibles publiquement, en jugeant que la notion
d'accès «~sans autorisation~» du Computer Fraud and Abuse Act ne s'applique pas aux sites publics. Dans l'Union européenne, le General Data
Protection Regulation (Regulation 2016/679) définit les données à caractère personnel comme toute information se rapportant à une personne physique
identifiée ou identifiable, directement ou indirectement, notamment par des données de localisation ou un identifiant en ligne.
\end{dated}

Une équipe de recherche tire de ces affaires une liste de contrôle~: la déclaration écrite du fournisseur sur la manière dont les données ont été collectées et sur les
consentements recueillis~; la question de savoir si une source est un initié ou est soumise à une obligation de confidentialité~; la présence éventuelle de données personnelles, pseudonymisées ou réellement agrégées~;
les conditions d'utilisation des sites dont les données ont été extraites (une injonction préliminaire portant sur des profils publics au titre d'une loi donnée n'est pas une licence d'extraction)~; un visa juridique et
de conformité avant l'essai, et non après l'achat~; et une trace écrite de l'ensemble, conservée avec les métadonnées du jeu de données. L'éthique va plus loin
que le droit~: des données permettant à un desk de déduire qui a fréquenté une clinique ne sont pas acceptables au motif qu'un tribunal ne s'est pas encore prononcé à leur sujet.

\section{La décision d'achat}

La valeur d'un jeu de données est le profit de ce qu'il apporte, net des coûts, sur la durée du contrat, à mesure que l'avantage s'érode parce que d'autres acheteurs
opèrent sur la même information. Dans l'essai, le spread long--court des quintiles supérieur et inférieur de la variable incrémentale rapporte 2.9~\% par trimestre pendant les
années en direct (variable brute en direct~: 3.4~\%~; les années reconstituées~: 10.4~\%). La capacité du desk pour les transactions sur annonces dans ces titres étant fixée à
20~millions de dollars, des coûts de transaction de 0.4~\% par an et une demi-vie de deux ans pour l'avantage, les frais annuels d'équilibre sur un contrat de trois ans s'élèvent à
1.35~million de dollars. L'historique du fournisseur, pris au pied de la lettre, implique 5.06~millions de dollars~; la variable brute en direct, 1.59~million de dollars
(\cref{fig:rs:alternative-and-text-data:breakeven}). La demi-vie tranche le reste~: à six mois, les frais d'équilibre incrémentaux sont de
0.43~million de dollars, proches du prix demandé.

\begin{omfigure}
\begin{tikzpicture}
\begin{axis}[width=10cm,height=5.6cm,xlabel={demi-vie de l'avantage (années)},ylabel={frais d'équilibre (millions de dollars par an)},
  tick label style={font=\small},label style={font=\small},xmin=0.4,xmax=5.1,ymin=0,ymax=7,grid=major,grid style={gray!25},
  legend columns=3,
  legend style={font=\footnotesize,at={(0.5,-0.28)},anchor=north,draw=none,fill=none,
  /tikz/every even column/.append style={column sep=8pt}},table/col sep=comma]
\addplot[gray,thick,dashed,mark=o] table[x=half_life,y=back]{figdata/research/12-alternative-and-text-data/breakeven.csv};
\addplot[omDef,thick,mark=*] table[x=half_life,y=live]{figdata/research/12-alternative-and-text-data/breakeven.csv};
\addplot[omThm,very thick,mark=square*] table[x=half_life,y=inc]{figdata/research/12-alternative-and-text-data/breakeven.csv};
\legend{historique reconstitué (discours commercial),années en direct,en direct et incrémental}
\end{axis}
\end{tikzpicture}

\omcaption{Frais annuels d'équilibre du panel synthétique de fiches sur un contrat de trois ans, en fonction de la demi-vie de son avantage, pour une capacité de 20~millions de dollars
et des coûts de transaction de 0.4~\% par an~: évalués sur l'historique reconstitué du fournisseur, sur les années en direct, et sur les années en direct nettes du
signal propre de la société. Données~: l'essai synthétique du chapitre.}\label{fig:rs:alternative-and-text-data:breakeven}
\end{omfigure}

Trois corrections séparent le discours commercial du prix~: données en direct seulement (le remplissage rétroactif a triplé le spread), incrémentales par rapport à ce qui est déjà détenu (un
supplémentaire de 15~\%), et érosion (un facteur trois entre une demi-vie de six mois et une demi-vie de deux ans). La demi-vie est la plus difficile à connaître avant
l'achat~: elle dépend du nombre d'autres acheteurs du même fichier, que le fournisseur connaît et que l'acheteur ignore~; le chapitre~28 la mesure après coup.
La capacité (chapitre~28) et les coûts (chapitre~27) proviennent des modèles propres de la société. Le problème du week-end du chapitre tranche la décision.

\section{Fiches de prédicteur}

\begin{predictorcard}{Surprise de ventes du panel de cartes}\label{pred:rs:alternative-and-text-data:card}
\sfield{Définition} La croissance des dépenses du panel pour le trimestre, résiduée sur les signaux existants de la société, classée entre les détaillants
qui annoncent dans les semaines à venir.
\sfield{Entrées et horodatages} Les valeurs trimestrielles du fournisseur avec leur date de première livraison (jamais la période)~; le calendrier des annonces
(chapitre~11).
\sfield{Justification} Des approximations des ventes intratrimestrielles expliquent les surprises et les rendements d'annonce (Froot, Kang, Ozik et Sadka).
\sfield{Horizon et demi-vie} De la livraison à l'annonce~; la demi-vie de l'avantage est fixée par le nombre d'acheteurs.
\sfield{Normalisation} Au sein de la coupe transversale du trimestre~; un indicateur de changement de panel comme variable de contrôle.
\sfield{Modes de défaillance} Historique reconstitué~; \omterm{def:rs:alternative-and-text-data:altdata}{dérive du panel}~; biais de couverture en faveur des commerçants du panel~; encombrement du même fichier par les acheteurs (crowding).
\sfield{Sources} Comme cité~; \texttt{rs\_altdata.trial\_report}.
\end{predictorcard}

\begin{predictorcard}{Ton selon le dictionnaire financier}\label{pred:rs:alternative-and-text-data:tone}
\sfield{Définition} Mots négatifs moins mots positifs, issus de listes propres à la finance, rapportés à l'ensemble des mots d'un dépôt, d'un communiqué ou d'un article de presse.
\sfield{Entrées et horodatages} L'heure de publication du document (l'heure d'acceptation du dépôt, l'horodatage de l'article)~; la version des
listes de mots.
\sfield{Justification} Le ton porte de l'information ou fait bouger les cours par le sentiment (Tetlock~; Loughran et McDonald).
\sfield{Horizon et demi-vie} Quelques jours après le document~; retournements longs pour les mouvements dictés par le sentiment.
\sfield{Normalisation} Par rapport aux documents passés de la même société et du même type (les formules types sont stables).
\sfield{Modes de défaillance} Listes d'usage général~; listes de mots ou modèles construits après la date de décision~; changements des formules types.
\sfield{Sources} Comme cité~; \texttt{rs\_altdata.text\_scores}.
\end{predictorcard}

\section{Tutoriel~: mener un essai de fournisseur}\label{tut:rs:alternative-and-text-data:tut}

\begin{tutorial}
\textbf{Objectif.} Lancer le banc d'essai sur le panel synthétique de fiches~: couverture, part de données reconstituées, lien du panel avec la vérité année par année et
sa rupture, IC par année, IC incrémental, spreads long--court et frais d'équilibre~; puis attribuer un score au corpus synthétique à l'aide de deux listes de mots.
\textbf{État final~:}
les \cref{fig:rs:alternative-and-text-data:drift,fig:rs:alternative-and-text-data:ic,fig:rs:alternative-and-text-data:breakeven}~; les
corrélations de texte.
\begin{tutsteps}
\item \textbf{Localiser la rupture}~: une CUSUM des résidus du panel par rapport à la vérité, dans l'ordre chronologique.
\omcode{code/firm/vendoreval/firm_vendoreval.py}{47}{53}{La rupture unique la plus vraisemblable dans la moyenne.}\label{lst:rs:alternative-and-text-data:cusum}
\item \textbf{IC incrémental}~: résidualiser au sein de chaque coupe transversale, puis calculer une corrélation de rangs.
\omcode{code/firm/vendoreval/firm_vendoreval.py}{69}{83}{Résidualisation sur les signaux détenus, et IC incrémental.}\label{lst:rs:alternative-and-text-data:inc}
\item \textbf{Chiffrer le prix}~: les frais qui égalent le profit net annuel moyen à mesure que l'avantage s'érode.
\omcode{code/firm/vendoreval/firm_vendoreval.py}{96}{101}{Les frais annuels d'équilibre.}\label{lst:rs:alternative-and-text-data:fee}
\item \textbf{Exécuter} \texttt{trial\_report()}, \texttt{breakeven\_curve()}, \texttt{text\_scores()} et \texttt{fig\_altdata.py}.
\end{tutsteps}
\textbf{Pour aller plus loin.} Masquer la date de lancement et tenter de la retrouver à partir des seules données~; donner au signal propre de la société davantage de
l'information du panel, et observer la baisse de l'IC incrémental et du prix.
\end{tutorial}

\section{Construction~: le harnais d'essai des fournisseurs}\label{bld:rs:alternative-and-text-data:vendoreval}

\begin{build}
\bfield{Rôle} Les mêmes questions posées à chaque jeu de données que la société envisage, avec des réponses comparables
d'un fournisseur à l'autre et versées au dossier de la décision d'achat.
\bfield{Interface} \texttt{coverage}, \texttt{backfill\_share(period, delivered, lag)}, \texttt{fit\_by\_group}, \texttt{cusum\_break},
\texttt{rank\_ic}, \texttt{ic\_by\_group}, \texttt{residualise}, \texttt{incremental\_ic(f, X, y, group)}, \texttt{long\_short},
\texttt{breakeven(annual, capital, cost, half\_life, years)}, \texttt{tone(counts, vocab, words)}.
\bfield{Règles} Aucun résultat d'essai sans dates de première livraison ni mention explicite de leur absence~; des IC par période, jamais seulement agrégés~;
une valeur calculée sur le prédicteur incrémental, nette des coûts et de l'érosion.
\bfield{Tests d'acceptation} \texttt{code/firm/vendoreval/tests/}~: couverture et part de remplissage rétroactif calculées à la main~; ajustements par groupe~; un décalage de moyenne injecté, retrouvé par
le CUSUM, et aucun dans le bruit~; un résidu non corrélé au signal détenu, un IC incrémental inférieur à l'IC brut et une copie sans rien de nouveau~; la
redevance sans érosion égale le bénéfice annuel~; parts de tonalité calculées à la main.
\bfield{Pour aller plus loin} Tests de rupture structurelle à plusieurs ruptures~; une mise à jour bayésienne de la demi-vie à partir des premiers mois de trading en direct~; une
liste de contrôle de la provenance stockée avec chaque jeu de données.
\end{build}

\begin{omsources}
\item K.~Froot, N.~Kang, G.~Ozik et R.~Sadka, ``What do measures of real-time corporate sales say about earnings surprises and post-announcement
returns?'', \emph{Journal of Financial Economics} 125(1), 2017.
\item C.~Zhu, ``Big data as a governance mechanism'', \emph{Review of Financial Studies} 32(5), 2019.
\item P.~C.~Tetlock, ``Giving content to investor sentiment: the role of media in the stock market'', \emph{Journal of Finance} 62(3), 2007.
\item T.~Loughran et B.~McDonald, ``When is a liability not a liability? Textual analysis, dictionaries, and 10-Ks'', \emph{Journal of Finance}
66(1), 2011.
\item US Securities and Exchange Commission, communiqué de presse 2021-176 (App Annie), 14 septembre 2021~; Staff Accounting Bulletin No.~99,
Materiality.
\item \emph{hiQ Labs, Inc.\ v.\ LinkedIn Corp.}, No.~17-16783 (9th Cir., 18 avril 2022).
\item Regulation (EU) 2016/679 (General Data Protection Regulation), article~4.
\end{omsources}

\section{Exercices}

\begin{exercise}[$\star$]\label{exo:rs:alternative-and-text-data:1}
Un fichier d'essai compte 1\,200 lignes pour les périodes antérieures au lancement du fournisseur et 1\,800 après. Sans information sur les dates de livraison, quelle part de
l'historique est exposée à un remplissage rétroactif, et que demanderiez-vous au fournisseur~?
\end{exercise}

\begin{exercise}[$\star$]\label{exo:rs:alternative-and-text-data:2}
Un document de 400 mots contient 12 mots d'une liste générale de termes négatifs, dont 9 sont «~liability~», «~tax~», «~cost~» ou «~capital~», et 5 mots positifs.
Calculer sa part de termes négatifs avec la liste générale et avec une liste qui exclut ces quatre mots, puis sa tonalité nette avec la seconde.
\end{exercise}

\begin{exercise}[$\star$]\label{exo:rs:alternative-and-text-data:3}
Une stratégie rapporte 2.9~\% par trimestre sur 20 millions de dollars, avec 0.4~\% de coûts par an et sans érosion. Quelle est la redevance annuelle au seuil de rentabilité~?
\end{exercise}

\begin{exercise}[$\star\star$]\label{exo:rs:alternative-and-text-data:4}
Le signal propre de la société et le panel du fournisseur mesurent tous deux la surprise $s$ avec des bruits indépendants~: $o = 0.45s + e_1$, $m = 0.35s + e_2$. Pourquoi
l'IC incrémental du fournisseur est-il positif mais inférieur à son IC brut~? Que se passe-t-il si $e_2$ est une copie de $e_1$~?
\end{exercise}

\begin{exercise}[$\star\star$]\label{exo:rs:alternative-and-text-data:5}
Montrer que la redevance au seuil de rentabilité d'un contrat de $Y$ années, dont le bénéfice annuel $A$ la première année est divisé par deux tous les $h$ ans (mesuré à mi-année), et sans
coûts, vaut $\frac{A}{Y}\sum_{k=0}^{Y-1} 2^{-(k+1/2)/h}$. L'évaluer pour $A = 2.3$ millions de dollars, $Y = 3$, $h = 2$.
\end{exercise}

\begin{exercise}[$\star\star$]\label{exo:rs:alternative-and-text-data:6}
Pourquoi un \omterm{def:rs:alternative-and-text-data:text}{plongement de document} issu d'un modèle de langage entraîné en 2025 fuit dans un backtest portant sur 2015 à 2020, même si chaque document est
correctement daté~?
\end{exercise}

\begin{exercise}[$\star\star\star$]\label{exo:rs:alternative-and-text-data:7}
\emph{Programmation.} Dans l'essai synthétique, masquer l'indicateur des lignes rétroactivement remplies et localiser la rupture à partir des IC du panel par trimestre, plutôt que de son ajustement à
la vérité. Quelle est la netteté de la rupture, et pourquoi l'ajustement à la vérité est-il le meilleur instrument~?
\end{exercise}

\begin{exercise}[$\star\star\star$]\label{exo:rs:alternative-and-text-data:8}
\emph{Trouver la faille.} «~Le fournisseur a confirmé que les données sont entièrement anonymisées, et la source est un site public, il n'y a donc aucune question juridique.~»
\end{exercise}

\section{Problème~: faut-il acheter le panel de cartes~?}

\begin{problem}[{Problème du week-end --- un essai de fournisseur, chiffré}]\label{pb:rs:alternative-and-text-data:1}
L'essai synthétique du chapitre~: 400 détaillants, cinq années d'un panel de cartes, un changement de fournisseur au début de la troisième année, un signal que la
société possède déjà, et un prix demandé de plusieurs centaines de milliers de dollars par an.

\medskip\noindent\textbf{Partie I --- Le fichier.}
\begin{enumerate}
\item Combien de lignes l'essai conserve-t-il, et pourquoi seulement les trimestres dont l'annonce est postérieure à la livraison~?
\item Combien de lignes par trimestre avant et après le lancement, et pourquoi la baisse est-elle un indice~?
\item Quelle part des lignes a été livrée plus d'un trimestre après sa période~?
\item Que doit fournir le fournisseur pour que cette part puisse être calculée~?
\end{enumerate}

\medskip\noindent\textbf{Partie II --- La dérive.}
\begin{enumerate}[resume]
\item Donner la pente, la corrélation et l'ordonnée à l'origine du panel par rapport à la vérité, chaque année.
\item Où le CUSUM place-t-il la rupture, et avec quelle statistique~?
\item Pourquoi la rupture coïncide-t-elle ici avec le changement de fournisseur, et quoi d'autre a changé à ce moment-là~?
\item Que conclure si le fournisseur n'avait pas mentionné le changement de fournisseur~?
\end{enumerate}

\medskip\noindent\textbf{Partie III --- L'information.}
\begin{enumerate}[resume]
\item Quels sont les IC par année, et l'IC du signal propre de la société pendant les années en direct~?
\item Quel est l'IC incrémental et sa statistique $t$~?
\item Quels sont les spreads long--short des années rétroactivement remplies, du prédicteur brut en direct et du prédicteur incrémental~?
\item Pourquoi la valeur doit-elle utiliser le prédicteur incrémental sur les seules données en direct~?
\end{enumerate}

\medskip\noindent\textbf{Partie IV --- La décision.}
\begin{enumerate}[resume]
\item Quelles redevances au seuil de rentabilité impliquent les trois spreads pour une demi-vie de deux ans~?
\item Quel est le seuil de rentabilité incrémental pour une demi-vie de six mois et de cinq ans~?
\item Énoncer le \emph{résultat nommé}~: le prix annuel au seuil de rentabilité du jeu de données, compte tenu de son IC incrémental, de sa capacité et de son érosion.
\item La société doit-elle acheter à 400\,000~\$ par an~? À 900\,000~\$~?
\item Que négocier, outre le prix~?
\item Quelles questions juridiques et de conformité doivent recevoir une réponse avant l'essai~?
\item Comment vérifier la demi-vie après l'achat~?
\item En une phrase~: que vaut un jeu de données~?
\end{enumerate}
\end{problem}

\section{Questions d'entretien}

\begin{interviewq}[$\star$ \iqroles{chercheur}]\label{iq:rs:alternative-and-text-data:1}
Citez trois façons dont l'historique d'un jeu de données alternatif peut fausser un backtest.
\end{interviewq}

\begin{interviewq}[$\star\star$ \iqroles{chercheur}]\label{iq:rs:alternative-and-text-data:2}
Comment décideriez-vous si un nouveau jeu de données apporte quelque chose aux signaux que vous possédez déjà~?
\end{interviewq}

\begin{interviewq}[$\star\star$ \iqroles{chercheur, ingénieur ML}]\label{iq:rs:alternative-and-text-data:3}
Vous notez des informations à l'aide d'un modèle de langage. Quels problèmes de \omterm{def:rs:point-in-time-data-and-the-biases:bitemporal}{temps de connaissance} devez-vous résoudre~?
\end{interviewq}

\begin{interviewq}[$\star\star$ \iqroles{chercheur, trader}]\label{iq:rs:alternative-and-text-data:4}
Un fournisseur demande 500\,000~\$ par an. Comment décidez-vous s'il faut payer~?
\end{interviewq}

\begin{interviewq}[$\star\star$ \iqroles{chercheur, risque}]\label{iq:rs:alternative-and-text-data:5}
Qu'est-ce qui fait des \omterm{def:rs:alternative-and-text-data:altdata}{données alternatives} un risque juridique, et quelle diligence mèneriez-vous~?
\end{interviewq}

\begin{interviewq}[$\star\star\star$ \iqroles{chercheur}]\label{iq:rs:alternative-and-text-data:6}
Pourquoi les dictionnaires de sentiment généralistes échouent-ils sur les textes financiers, et comment construiriez-vous une meilleure mesure~?
\end{interviewq}
